\section{Blue-team detection containment and recovery}\label{sec:C28:blueteam}
Blue-team controls operate at several boundaries. Input and retrieval controls preserve provenance and reduce untrusted instruction promotion. Tool controls enforce schemas, capabilities and principals. Runtime controls isolate processes, credentials and filesystems. Memory controls quarantine suspicious records and propagate corrections. Monitoring and recovery detect violations that prevention missed. These controls are complementary, but each has a cost and a specific failure surface.

Detection should be evaluated separately from prevention. A monitor can identify an unauthorized proposal that the tool layer already blocked, or miss an effect that occurred through an unmonitored channel. Report detection coverage, false alarms, latency and whether the intervention arrived before the effect. A high detection rate after irreversible disclosure is not equivalent to successful containment.

Quarantine protects future decisions while preserving evidence. A suspicious memory or skill can be excluded from active retrieval until reviewed. The system should retain provenance and dependencies so related records can be examined. Automatically deleting the evidence can prevent root-cause analysis; leaving it available to the model can continue the compromise. The quarantine namespace and review authority must therefore be application-controlled.

Recovery includes credential revocation, cancellation, state restoration and external-effect reconciliation. Restoring a clean checkpoint does not undo information already disclosed or actions already accepted by remote services. Report residual effects explicitly. A rollback that repairs local files but leaves a malicious scheduled task active has not restored the operational boundary.

The proposed experiment injects controlled failures at each layer and measures useful completion, violations, detection delay, containment and recovery. Include adaptive attacks against the defense configuration, not only payloads used to design it. Compare the full stack with single-control ablations under the same benign workload. This identifies where defense-in-depth supplies complementary protection and where several components merely repeat the same fallible classification.

\subsection{Detection must precede the relevant commitment}\label{sec:C28:precommitalert}
Let $T_A$ be the time a validated alert becomes actionable, $L_C\geq0$ the containment delay, and $T_E$ the counterfactual commitment time without that containment, under the declared fixture. A conservative timing requirement is $T_A+L_C<T_E$; an equality case depends on the service's ordering semantics. It is not sufficient: the containment action must control the actual effect channel and remain authorized. A prevented effect has no observed commitment timestamp, so its deadline must come from the fixture or a justified matched replay rather than be invented from the defender's narration. Reporting only alert recall discards these conditions. An asynchronous monitor can recognize every malicious trajectory and still prevent none if dispatch commits before the alert arrives.

Use service-side commitment receipts where available rather than the defender's message timestamp. For delayed or multi-stage effects, define the relevant point separately for each consequence: revoking access can stop future reads while leaving an earlier disclosure irreversible. Evaluate attempted violations, completed effects, alerts, successful containment and residual consequences as separate events. The measured prevention fraction is conditional on the fixture's commit and cancellation semantics. It cannot be inferred by applying a monitor score to a different tool whose effects commit synchronously or whose cancellation behavior is unavailable. This connects blue-team evaluation to the effect journal of Chapter~\ref{ch:C20} and the operational recovery analysis of Chapter~\ref{ch:C32}.
